% ===================================================================
%  TABLEAU N° 4 — L'âge mûr et la vieillesse.
%
%  Traduction, faite en 2026, en francais standard du Quebec, sur
%  l'ido de texto/io. Regles de la colonne : voir 10-tableau-01.tex.
%  Deux scenes, deux numerotations, comme au tableau 3 : les cles
%  portent c1 et c2.
%
%  LES HUIT CHOIX PROPRES A CE TABLEAU :
%    (25) docteur ............... médecin
%    (27) ordonnance ............ prescription
%    (48) bouteilles de liqueur . bouteilles d'alcool
%    (11) jaquette .............. veston
%    (14) bonne ................. domestique
%    (11) domestique (c1) ....... serveur
%    (23)(29) grand'maman, grand'mère  grand-maman, grand-mère
%         bisaïeul .............. arrière-grand-père
%
%  DEUX PIEGES, ET CE SONT LES DEUX MEMES FAUX AMIS DU QUEBEC.
%    * « liqueur » au Quebec est une boisson gazeuse. « Des bouteilles
%      de liqueur » sur une table de dessert y serait lu comme des
%      bouteilles de boisson gazeuse ; or l'ido dit « liquor-boteli »
%      et la gravure les pose a cote du champagne. C'est de l'alcool.
%    * « jaquette » au Quebec est une chemise de nuit. Le fac-simile
%      l'emploie pour le vetement de ville de la mere (11) : le garder
%      aurait mis la dame en tenue de nuit dans la rue, en plein
%      hiver.
%
%  (11) DU PREMIER TABLEAU ET (14) DU SECOND SONT DEUX PERSONNES
%  DIFFERENTES, et le fac-simile les nomme presque pareil : « le
%  domestique » qui dessert le repas de noces, « la bonne » qui porte
%  le bol de tisane. L'ido, lui, les distingue moins encore
%  (« servisto », « servistino »). On a donc tranche sur la scene :
%  celui du banquet fait le service d'un repas — c'est un serveur ;
%  celle de la maison est au service de la maison — c'est une
%  domestique. « Bonne » est parti pour son registre, non pour son
%  sens.
%
%  LES TROIS REPAS CHANGENT DE NOM, ET C'EST LE DECALAGE LE PLUS
%  REGULIER DE TOUTE LA COLONNE. Le Quebec dit dejeuner le matin,
%  diner le midi, souper le soir ; la France de 1926 disait dejeuner
%  le midi et diner le soir. Le repas de noces de ce tableau, qui se
%  prend « le soir » sous la veranda, est donc un SOUPER, et le
%  fac-simile l'appelle diner. La regle vaut pour toute la colonne :
%  « dineo » se rend par souper, « dejuno » par diner, chaque fois que
%  le texte situe le repas dans la journee. On la retrouvera aux
%  tableaux 5, 6 et 13.
%
%  LES DEUX PAIRES DE LUNETTES NE SONT PAS LES MEMES, et l'ido est
%  plus precis que le francais : « (naz)binoklo » (c1) est le
%  pince-nez du tuteur, « orelo-binoklo » (c2) sont les lunettes a
%  branches de la grand-mere. Le fac-simile dit « lunettes » pour la
%  seconde et « pince-nez » pour la premiere : on garde les deux.
% ===================================================================

%%K t04-apar-1 apar
\VUsaut{4.0mm}
\VUtitre{15.0pt}{60}{TABLEAU N\textsuperscript{o} 4}
\VUsaut{1.6mm}
\VUtitre{11.4pt}{0}{L'âge mûr et la vieillesse.}
\VUsaut{3.2mm}

%%K t04-tit-1 sub
\VUcentre{11.4pt}{0}{\textit{Première scène.}}
\VUsaut{1.6mm}
\VUtitre{11.4pt}{0}{Le repas de noces.}
\VUsaut{2.6mm}

%%K t04-tit-2 sub
\VUpk{10.2pt}{40}{L'automne.}
\VUsaut{3.0mm}

%%K t04-c1-01-1 p
1. --- Les jeunes ont grandi. L'un d'eux, Jean, se marie. Nous assistons au
\VUgras{repas de noces}, qui réunit autour de la même table les familles des
deux époux.

%%K t04-c1-02-1 p
2. --- Nous sommes en \VUgras{automne}, au mois de \VUgras{septembre}. Le vent
froid d'\VUgras{octobre} et de \VUgras{novembre} n'a pas encore dépouillé la
campagne de son feuillage vert. L'air du soir est tiède et parfumé; c'est
pourquoi le souper a lieu sous la \VUgras{véranda} \textsuperscript{(8)}, du
côté du jardin.

%%K t04-c1-03-1 p
3. --- Au loin, sur la \VUgras{colline} \textsuperscript{(3)}, se trouve la
maison des parents de Jean, un élégant \VUgras{château} \textsuperscript{(1)}
caché dans la verdure; de beaux vignobles, qui donnent un excellent vin,
couvrent le versant de la colline. Le vigneron les cultive et les soigne.

%%K t04-c1-04-1 p
4. --- Au-dessus des vignes passe un vol de \VUgras{perdrix}
\textsuperscript{(2)} effarouchées par l'approche du \VUgras{garde-chasse}
\textsuperscript{(5)}. Celui-ci, le \VUgras{fusil} sous le bras et la
\VUgras{gibecière} en bandoulière, bat les \VUgras{buissons}
\textsuperscript{(4)} avec son \VUgras{chien} \textsuperscript{(6)}. Il
surveille le domaine et empêche les braconniers de détruire le gibier.

%%K t04-c1-05-1 p
5. --- Au-dehors, contre la \VUgras{véranda}, grimpe une vigne en espalier à
laquelle pendent de lourdes \VUgras{grappes de raisin} \textsuperscript{(7)}.

%%K t04-c1-06-1 p
6. --- Les belles fleurs d'automne qui garnissent la \VUgras{table}
\textsuperscript{(16)} sont des \VUgras{chrysanthèmes} \textsuperscript{(9)};
le \VUgras{père} \textsuperscript{(10)} de la mariée en a mis une à la
boutonnière de son habit.

%%K t04-c1-07-1 p
7. --- Le repas touche à sa fin. Le \VUgras{serveur} \textsuperscript{(11)}
commence à apporter les plats de dessert et enlèvera bientôt les différentes
pièces du couvert, \VUgras{cuillères} \textsuperscript{(14)},
\VUgras{fourchettes} \textsuperscript{(12)}, \VUgras{couteaux}
\textsuperscript{(13)} et \VUgras{plats} \textsuperscript{(15)}, dont on n'a
plus besoin.

%%K t04-tit-3 sub
\VUpk{10.2pt}{40}{La parenté.}
\VUsaut{3.0mm}

%%K t04-c1-08-1 p
8. --- Tout le monde a l'air joyeux, et le sourire avec lequel la
\VUgras{mère} \textsuperscript{(17)} du marié regarde ses enfants prouve son
bonheur.

%%K t04-c1-09-1 p
9. --- Ce mariage unit deux riches familles du pays. Jean, le \VUgras{mari}
\textsuperscript{(18)}, est le \VUgras{fils} d'un grand propriétaire et il
devient le \VUgras{gendre} d'un habile industriel.

%%K t04-c1-10-1 p
10. --- Les époux sont jeunes, et pourtant ils étaient \VUgras{fiancés} depuis
longtemps. La \VUgras{jeune femme} \textsuperscript{(19)}, Marthe, est
charmante dans sa \VUgras{robe blanche} garnie de fleurs d'oranger.

%%K t04-c1-11-1 p
11. --- Son \VUgras{beau-père} \textsuperscript{(20)} est très heureux d'avoir
une \VUgras{bru} aussi gentille, et la \VUgras{belle-mère}
\textsuperscript{(21)} de Jean sourit en voyant sa \VUgras{fille} si belle.

%%K t04-c1-12-1 p
12. --- Au bout de la table sont assis les autres \VUgras{invités}
\textsuperscript{(22)} : \VUgras{parents, amis, cousins, cousines}. Un
\VUgras{maître d'hôtel} \textsuperscript{(23)}, la serviette sous le bras, les
sert avec empressement.

%%K t04-tit-4 sub
\VUpk{10.2pt}{40}{Les amis.}
\VUsaut{3.0mm}

%%K t04-c1-13-1 p
13. --- Du côté droit de la table, voici une \VUgras{demoiselle}
\textsuperscript{(24)}, \VUgras{amie} de la jeune mariée, puis le
\VUgras{frère} de celle-ci, qui est son \VUgras{garçon d'honneur}
\textsuperscript{(25)}. Il bavarde avec sa \VUgras{demoiselle d'honneur}
\textsuperscript{(26)}, la sœur du marié, qui devient par ce mariage la
\VUgras{belle-sœur} de Marthe. C'est une jeune fille charmante. Elle a de beaux
\VUgras{bijoux}, un \VUgras{collier de perles} et un \VUgras{bracelet} d'or.

%%K t04-c1-14-1 p
14. --- Près d'eux, le monsieur qui se tient debout et lève son verre est un
\VUgras{ami} \textsuperscript{(27)} de la famille. Il est l'un des
\VUgras{témoins du marié}. Il porte un \VUgras{toast}, il boit à la
\VUgras{santé} des jeunes mariés et leur souhaite un long bonheur. Il est en
tenue de cérémonie : habit et \VUgras{souliers vernis} \textsuperscript{(28)}.
C'est lui qui accompagne la \VUgras{grand-mère} \textsuperscript{(29)}, qui est
\VUgras{veuve}.

%%K t04-c1-15-1 p
15. --- Le jeune homme en \VUgras{habit} \textsuperscript{(31)}, à la fine
moustache noire, est le \VUgras{beau-frère} \textsuperscript{(30)} de Jean.
C'est un ingénieur remarquable. Il est le voisin d'une \VUgras{jeune femme}
\textsuperscript{(33)} très élégante, la \VUgras{sœur aînée} de Marthe et la
mère de la mignonne petite fille qui vient, au bras de son cousin, lui offrir
une fleur et lui \VUgras{souhaiter} le \VUgras{bonsoir}. Cette dame a de très
belles \VUgras{boucles d'oreilles} et une élégante \VUgras{coiffure}.

%%K t04-c1-16-1 p
16. --- Le petit cousin, si drôle avec sa \VUgras{blouse} \textsuperscript{(36)}
et sa \VUgras{culotte} \textsuperscript{(37)} bouffante, est bien à plaindre.
Il est \VUgras{orphelin} \textsuperscript{(35)}. Il a perdu tout jeune son père
et sa mère. Son oncle lui sert de \VUgras{tuteur} \textsuperscript{(38)} et est
resté célibataire pour élever le pauvre enfant. C'est un homme bon. Il est un
peu chauve et très myope; il porte un \VUgras{pince-nez}.

%%K t04-tit-5 sub
\VUsaut{2.4mm}
\VUpk{10.2pt}{40}{Le dessert.}
\VUsaut{3.0mm}

%%K t04-c1-17-1 p
17. --- Derrière les convives, sur une table recouverte d'une \VUgras{nappe},
le \VUgras{dessert} \textsuperscript{(39)} est préparé. Dans une coupe, voici
des fruits : des \VUgras{pêches} \textsuperscript{(40)}, des \VUgras{poires}
\textsuperscript{(41)}, des \VUgras{pommes} \textsuperscript{(42)}, des
\VUgras{abricots} \textsuperscript{(43)} et des \VUgras{raisins}
\textsuperscript{(44)}. À côté, des \VUgras{ananas} \textsuperscript{(45)}, un
beau \VUgras{gâteau} \textsuperscript{(46)}, des \VUgras{bouteilles d'alcool}
\textsuperscript{(48)} et du \VUgras{champagne} \textsuperscript{(49)}, ainsi
que des piles d'\VUgras{assiettes} \textsuperscript{(47)} propres.

%%K t04-c1-18-1 p
18. --- Le \VUgras{sommelier} \textsuperscript{(50)} débouche une
\VUgras{bouteille} \textsuperscript{(52)} de vin vieux avec un
\VUgras{tire-bouchon} \textsuperscript{(51)}. Le \VUgras{bouchon}
\textsuperscript{(53)} de liège a l'air bien difficile à tirer. Ce sommelier a
une \VUgras{serviette} \textsuperscript{(54)} sous le bras.

%%K t04-c1-19-1 p
19. --- Après le souper, les messieurs iront au fumoir, et les dames iront se
préparer pour le bal.

%%K t04-tit-6 sub
\VUsaut{5.0mm}
\VUcentre{11.4pt}{0}{\textit{Deuxième scène.}}
\VUsaut{1.6mm}
\VUpk{11.4pt}{40}{L'anniversaire du grand-père.}
\VUsaut{2.6mm}

%%K t04-tit-7 sub
\VUpk{10.2pt}{40}{La visite.}
\VUsaut{3.0mm}

%%K t04-c2-01-1 p
1. --- Dix ans ont passé. Les parents de Jean sont devenus vieux et aiment
beaucoup leurs trois petits-enfants, deux \VUgras{petits-fils}
\textsuperscript{(2)} et une \VUgras{petite-fille} \textsuperscript{(3)}. Comme
c'est aujourd'hui l'\VUgras{anniversaire du grand-père}, les enfants viennent
lui souhaiter une bonne fête.

%%K t04-c2-02-1 p
2. --- Le petit Pierre est encore très jeune, il n'a que trois ans. Il voit
s'approcher le \VUgras{chat} \textsuperscript{(1)}. Il veut caresser son beau
\VUgras{pelage} soyeux et sa longue \VUgras{queue}; il ne songe pas que sous
ses gracieuses \VUgras{pattes} se cachent de méchantes griffes pointues qui
risquent de l'égratigner.

%%K t04-c2-03-1 p
3. --- Sa sœur Gabrielle, qui lui donne la main, est bien plus sérieuse, et
Marcel, avec son grand \VUgras{manteau} \textsuperscript{(4)} et sa
\VUgras{ceinture} \textsuperscript{(5)} de cuir, a l'air d'un petit homme,
malgré ses culottes courtes qui laissent voir ses \VUgras{bas}
\textsuperscript{(6)}.

%%K t04-c2-04-1 p
4. --- Leur père a mis son gros \VUgras{pardessus} \textsuperscript{(7)}
d'hiver à \VUgras{col de fourrure} \textsuperscript{(8)} et il a son foulard
dans sa \VUgras{poche} \textsuperscript{(9)}. Sa femme a son chapeau de feutre
garni de \VUgras{plumes} \textsuperscript{(10)}, son \VUgras{veston}
\textsuperscript{(11)}, son \VUgras{boa} \textsuperscript{(12)} et son
\VUgras{manchon} \textsuperscript{(13)}.

%%K t04-c2-05-1 p
5. --- Il fait très froid, parce que c'est l'hiver. La \VUgras{neige}
\textsuperscript{(19)} est tombée toute la journée et les toits des maisons
sont tout blancs. Avec la \VUgras{nuit}, le froid devient encore plus vif. Dans
le ciel, la \VUgras{lune} \textsuperscript{(17)} et les \VUgras{étoiles}
\textsuperscript{(18)} brillent et percent l'\VUgras{obscurité}.

%%K t04-tit-8 sub
\VUsaut{4.2mm}
\VUpk{10.2pt}{40}{La maladie.}
\VUsaut{3.0mm}

%%K t04-c2-06-1 p
6. --- Le pauvre \VUgras{aveugle} \textsuperscript{(20)} qui traverse la place
devant la \VUgras{pharmacie} \textsuperscript{(16)}, conduit par son
\VUgras{caniche} \textsuperscript{(21)}, est bien malheureux. Justine, la
\VUgras{domestique} \textsuperscript{(14)}, apporte de la cuisine un
\VUgras{bol} \textsuperscript{(15)} de \VUgras{tisane} bien chaude et
généreusement sucrée.

%%K t04-c2-07-1 p
7. --- La \VUgras{grand-maman} \textsuperscript{(23)}, qui veut aller chercher
ses \VUgras{lunettes} \textsuperscript{(26)} sur la table pour lire la
\VUgras{prescription} \textsuperscript{(27)} que le \VUgras{médecin}
\textsuperscript{(25)} vient d'écrire, se lève de son fauteuil en s'appuyant
sur sa \VUgras{canne} \textsuperscript{(24)}.

%%K t04-c2-08-1 p
8. --- Elle souffre souvent de petits \VUgras{malaises}, de \VUgras{vertiges},
de \VUgras{rhumes} et de \VUgras{maux de dents}; ses jambes sont faibles et ses
yeux ne sont plus très bons. Malgré tout, elle se moque volontiers de son vieux
\VUgras{médecin}, qui veut toujours lui prescrire une \VUgras{potion}
\textsuperscript{(28)}. Aujourd'hui, elle est très joyeuse d'avoir autour
d'elle sa jeune famille.

%%K t04-c2-09-1 p
9. --- On est bien dans la chambre, portes closes. Dans la \VUgras{cheminée}
\textsuperscript{(29)} de marbre flambe un \VUgras{beau feu}
\textsuperscript{(30)}, et l'on se sent heureux sous la douce \VUgras{lumière}
de la \VUgras{lampe} \textsuperscript{(31)} qu'on vient d'allumer.

%%K t04-tit-9 sub
\VUsaut{4.2mm}
\VUpk{10.2pt}{40}{Le grand-père.}
\VUsaut{3.0mm}

%%K t04-c2-10-1 p
10. --- Le \VUgras{grand-père} \textsuperscript{(33)} lui-même a oublié qu'il
était malade, que son \VUgras{rhumatisme} le cloue dans son \VUgras{fauteuil}
\textsuperscript{(34)} et qu'hier encore il avait de la \VUgras{fièvre} et des
\VUgras{syncopes}. Il ne se rappelle plus que l'hiver dernier il a souffert
d'une \VUgras{pneumonie} et qu'une \VUgras{toux} rauque et pénible l'a fait
beaucoup souffrir.

%%K t04-c2-10-2 p
Et pourtant il a eu bien du mal à guérir. Depuis, il va un peu mieux. Mais sa
jambe, qu'il tient appuyée sur un \VUgras{coussin} \textsuperscript{(38)} posé
sur un petit \VUgras{tabouret} \textsuperscript{(39)}, le fait souffrir elle
aussi.

%%K t04-c2-11-1 p
11. --- Quand il est soigneusement enveloppé dans sa chaude \VUgras{robe de
chambre} \textsuperscript{(35)} à gros \VUgras{boutons}
\textsuperscript{(36)} de nacre, et qu'il a près de lui, pour lui faire la
lecture du journal, la petite \VUgras{orpheline} \textsuperscript{(40)} en
\VUgras{coiffe} blanche, en \VUgras{robe} \textsuperscript{(42)} bleue et en
\VUgras{pèlerine} \textsuperscript{(41)}, il ne se plaint pas.

%%K t04-c2-12-1 p
12. --- Chaque année, lorsque le \VUgras{calendrier} \textsuperscript{(43)}
marque de nouveau la \VUgras{date} du premier \VUgras{décembre}, il attend
impatiemment ses \VUgras{petits-enfants}, car il sait qu'ils n'oublieront pas
cette \VUgras{date} importante.

%%K t04-c2-13-1 p
13. --- Il se rappelle avec émotion le temps, bien lointain, où il venait
lui-même souhaiter bonne fête au glorieux \VUgras{maréchal} d'Empire dont le
\VUgras{portrait} \textsuperscript{(44)} est accroché au mur, et qui est
l'\VUgras{arrière-grand-père} de ses enfants.
\VUsaut{6.0mm}
\VUfilet{22mm}
